% GENERATED FILE. Edit canonical lessons, then run npm run generate:books.
% canonical-source-hash: fnv1a64:2aad37b629dd35db
% canonical-lessons: KA-C132-navilu, KA-C132-citte, KA-C132-iruve, KA-C132-solle, KA-C132-jeda

\chapter{Peacock, Butterfly, Ant, Mosquito, Spider}
\label{ch:ka-peacock-butterfly-ant-mosquito-spider}
\hlchaptermodality{132}

\begin{chapteropening}
\textbf{By the end of this chapter:} \emph{I can say a peacock, a butterfly, an ant, a mosquito and a spider.}

Complete the last lesson of chapter 132: I can say a peacock, a butterfly, an ant, a mosquito and a spider.
\end{chapteropening}

\section[\texorpdfstring{\kn{ನವಿಲು} (navilu)}{ನವಿಲು (navilu)}]{\kn{ನವಿಲು} (navilu) — a peacock}

\label{lesson:KA-C132-navilu}



\begin{quote}
Before the new one: say the Kannada for a pigeon, then the Kannada for an owl.
\end{quote}

\subsection*{You'll want to know: \kn{ನವಿಲು}}
\textbf{\kn{ನವಿಲು}} — \emph{navilu} — \textquotedblleft{}a peacock\textquotedblright{}.

\begin{grammarlens}[title={What you've built}]
One more word for this chapter.
\end{grammarlens}

\subsection*{Your turn}
\begin{itemize}
\raggedright
  \item \emph{Say it:} \emph{navilu}
  \item \emph{Say it:} \emph{navilu}, once more
  \item \emph{Say it:} say \emph{navilu}
  \item \emph{Recall:} say the Kannada for a pigeon, then the Kannada for an owl, then say \emph{navilu} again
\end{itemize}

\begin{tcolorbox}[breakable,colback=teal!4,colframe=teal!35!black,title={Before you move on}]
What does \kn{ನವಿಲು} mean? (\textquotedblleft{}A peacock\textquotedblright{}.) Say it once more.
\end{tcolorbox}
\section[\texorpdfstring{\kn{ಚಿಟ್ಟೆ} (ciṭṭe)}{ಚಿಟ್ಟೆ (ciṭṭe)}]{\kn{ಚಿಟ್ಟೆ} (ciṭṭe) — a butterfly}

\label{lesson:KA-C132-citte}



\begin{quote}
Before the new one: say the Kannada for an owl, then the Kannada for a peacock.
\end{quote}

\subsection*{You'll want to know: \kn{ಚಿಟ್ಟೆ}}
\textbf{\kn{ಚಿಟ್ಟೆ}} — \emph{ciṭṭe} — \textquotedblleft{}a butterfly\textquotedblright{}.

\begin{grammarlens}[title={What you've built}]
One more word for this chapter.
\end{grammarlens}

\subsection*{Your turn}
\begin{itemize}
\raggedright
  \item \emph{Say it:} \emph{ciṭṭe}
  \item \emph{Say it:} \emph{ciṭṭe}, once more
  \item \emph{Say it:} say \emph{ciṭṭe}
  \item \emph{Recall:} say the Kannada for an owl, then the Kannada for a peacock, then say \emph{ciṭṭe} again
\end{itemize}

\begin{tcolorbox}[breakable,colback=teal!4,colframe=teal!35!black,title={Before you move on}]
What does \kn{ಚಿಟ್ಟೆ} mean? (\textquotedblleft{}A butterfly\textquotedblright{}.) Say it once more.
\end{tcolorbox}
\section[\texorpdfstring{\kn{ಇರುವೆ} (iruve)}{ಇರುವೆ (iruve)}]{\kn{ಇರುವೆ} (iruve) — an ant}

\label{lesson:KA-C132-iruve}



\begin{quote}
Before the new one: say the Kannada for a peacock, then the Kannada for a butterfly.
\end{quote}

\subsection*{You'll want to know: \kn{ಇರುವೆ}}
\textbf{\kn{ಇರುವೆ}} — \emph{iruve} — \textquotedblleft{}an ant\textquotedblright{}.

\begin{grammarlens}[title={What you've built}]
One more word for this chapter.
\end{grammarlens}

\subsection*{Your turn}
\begin{itemize}
\raggedright
  \item \emph{Say it:} \emph{iruve}
  \item \emph{Say it:} \emph{iruve}, once more
  \item \emph{Say it:} say \emph{iruve}
  \item \emph{Recall:} say the Kannada for a peacock, then the Kannada for a butterfly, then say \emph{iruve} again
\end{itemize}

\begin{tcolorbox}[breakable,colback=teal!4,colframe=teal!35!black,title={Before you move on}]
What does \kn{ಇರುವೆ} mean? (\textquotedblleft{}An ant\textquotedblright{}.) Say it once more.
\end{tcolorbox}
\section[\texorpdfstring{\kn{ಸೊಳ್ಳೆ} (soḷḷe)}{ಸೊಳ್ಳೆ (soḷḷe)}]{\kn{ಸೊಳ್ಳೆ} (soḷḷe) — a mosquito}

\label{lesson:KA-C132-solle}



\begin{quote}
Before the new one: say the Kannada for a butterfly, then the Kannada for an ant.
\end{quote}

\subsection*{You'll want to know: \kn{ಸೊಳ್ಳೆ}}
\textbf{\kn{ಸೊಳ್ಳೆ}} — \emph{soḷḷe} — \textquotedblleft{}a mosquito\textquotedblright{}.

\begin{grammarlens}[title={What you've built}]
One more word for this chapter.
\end{grammarlens}

\subsection*{Your turn}
\begin{itemize}
\raggedright
  \item \emph{Say it:} \emph{soḷḷe}
  \item \emph{Say it:} \emph{soḷḷe}, once more
  \item \emph{Say it:} say \emph{soḷḷe}
  \item \emph{Recall:} say the Kannada for a butterfly, then the Kannada for an ant, then say \emph{soḷḷe} again
\end{itemize}

\begin{tcolorbox}[breakable,colback=teal!4,colframe=teal!35!black,title={Before you move on}]
What does \kn{ಸೊಳ್ಳೆ} mean? (\textquotedblleft{}A mosquito\textquotedblright{}.) Say it once more.
\end{tcolorbox}
\section[\texorpdfstring{\kn{ಜೇಡ} (jēḍa)}{ಜೇಡ (jēḍa)}]{\kn{ಜೇಡ} (jēḍa) — a spider}

\label{lesson:KA-C132-jeda}



\begin{quote}
Before the new one: say the Kannada for an ant, then the Kannada for a mosquito.
\end{quote}

\subsection*{You'll want to know: \kn{ಜೇಡ}}
\textbf{\kn{ಜೇಡ}} — \emph{jēḍa} — \textquotedblleft{}a spider\textquotedblright{}.

\begin{grammarlens}[title={What you've built}]
One more word for this chapter.
\end{grammarlens}

\subsection*{Your turn}
\begin{itemize}
\raggedright
  \item \emph{Say it:} \emph{jēḍa}
  \item \emph{Say it:} \emph{jēḍa}, once more
  \item \emph{Say it:} say \emph{jēḍa}
  \item \emph{Recall:} say the Kannada for an ant, then the Kannada for a mosquito, then say \emph{jēḍa} again
\end{itemize}

\begin{tcolorbox}[breakable,colback=teal!4,colframe=teal!35!black,title={Before you move on}]
What does \kn{ಜೇಡ} mean? (\textquotedblleft{}A spider\textquotedblright{}.) Say it once more.
\end{tcolorbox}
